\chapter{Causal Machine Learning}

\section{From prediction to treatment effects}

Machine-learning estimators can also be used for causal inference when
the goal is to estimate a treatment effect $Y(1) - Y(0)$. The commands in
this chapter target two quantities:

\begin{itemize}
\item the \textbf{Average Treatment Effect} (ATE)
\item the \textbf{Conditional Average Treatment Effect} (CATE),
      $\tau(x) = \E[Y(1) - Y(0) \mid X = x]$
\end{itemize}

They all require a treatment indicator \texttt{d} and a set of covariates \texttt{x}.
They do not replace a clear research design, but they can help discover
heterogeneous effects and improve precision.

\section{Causal Forest}

Causal Forest (Athey and Wager, 2018) adapts Random Forest to estimate
heterogeneous treatment effects directly. It splits on covariates to
maximise the variance of estimated treatment effects across leaves.

\begin{lstlisting}[caption={Causal Forest}]
set_seed(42)
let n = 1000
generate df x1 = rnormal(n)
generate df x2 = rnormal(n)
generate df d = (0.5 + 0.3*x1 > runiform(n))
generate df y = 1 + 2*x1 - x2 + d*(0.5 + 0.4*x1) + rnormal(n)

let m_cf = causalforest(y ~ d, df, x="x1,x2", trees=200, depth=5)
print(m_cf)
\end{lstlisting}

The output includes:

\begin{itemize}
\item the estimated ATE
\item a CATE distribution (mean, quartiles)
\item variable importance
\end{itemize}

\section{DR-Learner and TMLE}

The \textbf{DR-Learner} (Kennedy, 2023) and \textbf{TMLE} (van der Laan
and Rubin, 2006) are doubly-robust approaches. Each constructs a pseudo-
outcome from estimated nuisance functions and then regresses it on
covariates to recover the CATE or ATE.

\begin{lstlisting}[caption={DR-Learner}]
let m_dr = dr_learner(y ~ d, df, x="x1,x2", folds=5)
print(m_dr)
\end{lstlisting}

\begin{lstlisting}[caption={TMLE}]
let m_tmle = tmle(y ~ d, df, w="x1,x2")
print(m_tmle)
\end{lstlisting}

Doubly-robust methods are attractive because the ATE remains consistent
if either the propensity score or the outcome model is well specified.

\section{Orthogonal Random Forest}

ORF (Oprescu, Syrgkanis, and Wu, 2019) residualises treatment and
outcome with respect to confounders before running a forest on the
residualised signals.

\begin{lstlisting}[caption={Orthogonal Random Forest}]
let m_orf = orf(y ~ d, df, x="x1,x2", w="x1,x2", trees=100, depth=5)
print(m_orf)
\end{lstlisting}

It is useful when the treatment is not randomised and the relevant
heterogeneity comes from a subset of covariates.

\section{Bayesian Synthetic Control}

\texttt{bsc} implements a Bayesian version of the synthetic-control estimator.
It is appropriate when a single treated unit is compared against several
donor units and the goal is a counterfactual path with credible
intervals.

\begin{lstlisting}[caption={Bayesian synthetic control}]
// df has columns y, c1, c2, ..., t
let m_bsc = bsc(df, y, "c1,c2", 25, prior=1.0)
print(m_bsc)
\end{lstlisting}

The output reports the treatment effect $\tau$, its standard deviation,
a 95\% credible interval, and the synthetic-control weights.

\section{Double Machine Learning}

\texttt{dml\_crossfit} estimates the ATE in a partially linear model

\[
Y = \theta D + g(X) + \varepsilon
\]

using K-fold cross-fitting for the nuisance function $g(X)$.

\begin{lstlisting}[caption={Double Machine Learning}]
let m_dml = dml_crossfit(y ~ d, df, x="x1,x2", folds=5)
print(m_dml)
\end{lstlisting}

Because the nuisance is estimated on a separate fold, the resulting ATE
has a standard $\sqrt{n}$-rate distribution and honest standard errors.

\section{Choosing a method}

\begin{itemize}
\item \textbf{Heterogeneous effects:} Causal Forest or ORF
\item \textbf{Doubly-robust ATE:} DR-Learner, TMLE, DML
\item \textbf{Single treated unit:} BSC
\end{itemize}

All of these commands are \textit{experimental}. Validation against
reference implementations is partial or not supported for some of them.
Use them for exploration and as complements to a clear design, not as
substitutes for it.
